% Generated draft: requires mathematical and source review.
\documentclass[11pt]{article}
\usepackage[margin=25mm]{geometry}
\usepackage{fontspec}
\setmainfont{Latin Modern Roman}
\usepackage{amsmath,amssymb,mathtools,unicode-math}
\setmathfont{Latin Modern Math}
\usepackage{enumitem,needspace,xurl,hyperref}
\setlength{\parindent}{0pt}
\setlength{\parskip}{5pt}
\setlength{\emergencystretch}{3em}
\newcommand{\N}{\mathbb{N}}
\newcommand{\Z}{\mathbb{Z}}
\newcommand{\Q}{\mathbb{Q}}
\newcommand{\R}{\mathbb{R}}
\newcommand{\C}{\mathbb{C}}
\newcommand{\F}{\mathbb{F}}
\newcommand{\A}{\mathbb{A}}
\newcommand{\PP}{\mathbb{P}}
\newcommand{\E}{\mathbb{E}}
\newcommand{\eps}{\varepsilon}
\newcommand{\dd}{\,\mathrm{d}}
\begin{document}
\section*{UnsolvedMath ID: 30001515}
\textbf{Source problem code:} catalogue entry
\subsection{Definitions and mathematical statement}
Let \((N,g)\) be a Riemannian manifold and let \(M\) be an \(n\)-dimensional area-minimizing hypersurface in \(N\). Write \(B_r^g(x)\) for the ambient geodesic ball and \(\omega_n\) for the Euclidean volume of the unit ball in \(\R^n\). With \(\mu_M\) denoting the area measure associated with the chosen minimizing object, the usual normalized density is
\[
\Theta(M,x)=\lim_{r\downarrow0}\frac{\mu_M\bigl(B_r^g(x)\bigr)}{\omega_n r^n}.
\]
If the minimizing object carries multiplicity, its mass measure counts that multiplicity. Let \(\operatorname{Sing}_{\mathrm{int}}(M)\) denote the singular points of \(M\) in the ambient interior, away from any prescribed boundary; regularity is understood in the chosen category of minimizing hypersurfaces.

The supplied question asks for the value of
\[
\inf\left\{\Theta(M,x):\ (N,g)\text{ is a Riemannian manifold},\quad M\subset N\text{ is an area-minimizing hypersurface},\quad x\in\operatorname{Sing}_{\mathrm{int}}(M)\right\}.
\]
The quantifier ranges over all such ambient manifolds, hypersurfaces, and interior singular points; the source imposes no fixed dimension. However, it does not specify the category of minimizing objects, allowed multiplicities or coefficients, or the precise local or global meaning of area minimization. These choices remain unresolved qualifications of the displayed infimum. The density formula above states the usual normalization, which is not explicitly defined in the supplied source.
\subsection{Short English statement}
Determine the infimum of the normalized area density at interior singular points, ranging over all area-minimizing hypersurfaces in Riemannian manifolds. The intended minimizing category, multiplicity conventions, and dimension conventions require clarification.
\subsection{Sources}
\begin{itemize}
\item Supplied local UnsolvedMath registry record: permanent ID \(30001515\), source problem code \texttt{catalogue entry}, title ``Minimum Density of Area-Minimizing Hypersurface Singularities,'' in the Oberwolfach Reports Open Problems collection. No individual proposer is identified; the supplied proposed year is 2010.
\item Supplied report citation: \emph{Calculus of Variations (2011)}, Oberwolfach Reports, pp.\ 1827--1884, DOI \texttt{10.4171/owr/2010/31}. \url{https://doi.org/10.4171/owr/2010/31}. The background attributes the problem to workshop open-problem discussions documented in Oberwolfach Reports.
\item Supplied literature reference: Ilmanen--White, \emph{Sharp lower bounds on density for area-minimizing cones}. \url{https://doi.org/10.4310/cjm.2015.v3.n1.a1}.
\item Supplied literature reference labeled ``Original Oberwolfach problem report'': \url{https://doi.org/10.4171/owr/2010/31}.
\item Literature assessment dated exactly \texttt{2026-08-21}: supplied dated curation, not independently verified. It labels the problem partially solved, asserts equivalence with the isolated-singularity cone problem, and attributes a sharp lower bound for topologically nontrivial minimizing hypercones to Ilmanen and White. It explicitly states that this theorem does not settle the unrestricted infimum over all minimizing hypersurfaces. The supplied background describes this assessment as literature triage derived from OWR curation and calls for independent verification of time-sensitive status and cited claims.
\item The registry labels the statement \texttt{corrected\_verified}; its accompanying explanation says that the complete question was preserved while mathematical notation was placed in delimiters. This supplied label does not constitute independent mathematical verification.
\end{itemize}
\end{document}
